% ============================================================
%   Chapter 3  \textemdash\  Requirements Analysis
% ============================================================

\chapter{Requirements Analysis}\label{ch:requirements}

This chapter captures the requirements that guided the design of the framework. It distinguishes
functional requirements (what the system has to \emph{do}, Section~\ref{sec:req-functional}) from
non-functional requirements (how \emph{well} it has to do it, Section~\ref{sec:req-nonfunctional}).
Section~\ref{sec:req-constraints} lists the technical constraints imposed by the simulator, the
hardware and the academic timeline of the project. Section~\ref{sec:req-acceptance} converts
these requirements into measurable acceptance criteria that can be evaluated with the metrics
logged by the training pipeline.

\section{Functional requirements}\label{sec:req-functional}

Functional requirements describe the observable behaviour of the system. Each requirement is
given a short identifier of the form \texttt{FR-XX}, a descriptive title and a succinct
rationale. Requirements that depend on a specific implementation choice are justified in
Chapter~\ref{ch:design}.

\begin{table}[H]
  \centering
  \caption{Functional requirements of CALA-SafeRL.}
  \label{tab:functional-requirements}
  \small
  \begin{tabular}{@{}llp{7.8cm}@{}}
    \toprule
    \textbf{ID} & \textbf{Title} & \textbf{Description} \\
    \midrule
    FR-01 & Gymnasium environment    & The simulator shall be exposed through a
        \texttt{gymnasium.Env}-compatible wrapper supporting \texttt{reset}, \texttt{step},
        \texttt{render} and \texttt{close}. \\
    FR-02 & Continuous control       & The action space shall be a two-dimensional continuous
        box \([-1,1]^2\) corresponding to steering and joint throttle\textendash brake. \\
    FR-03 & Semantic LIDAR           & The observation shall include three LIDAR channels
        (combined, dynamic, static) of 240 rays each, normalised to \([0,1]\). \\
    FR-04 & Lane and route features  & The observation shall include lane features (lateral
        offset, heading error, edge distances, curvature, lane-change direction) and route
        features (waypoint angles, progress, speed limit) derived from the Waypoint API. \\
    FR-05 & PPO training loop        & A PPO agent with GAE, clipped surrogate, value loss,
        entropy bonus and KL early stopping shall be provided. \\
    FR-06 & Interchangeable shields  & The training and evaluation entry points shall accept a
        \texttt{--shield\_type} argument with values \texttt{none}, \texttt{basic} or
        \texttt{adaptive}, and switch shield implementations without code changes. \\
    FR-07 & Dense reward shaping     & A wrapper shall augment CARLA's sparse reward with dense
        components (speed, lane centring, heading, smoothness, progress, idle, shield
        interaction). \\
    FR-08 & Performance-driven curriculum & A curriculum manager shall progressively add NPC
        traffic as the agent's performance improves and shall roll back when performance
        collapses. \\
    FR-09 & Checkpoint persistence   & The agent shall save the best checkpoint on reward
        improvement and periodic checkpoints at a configurable frequency. \\
    FR-10 & Live metrics             & Every episode and every PPO update shall be logged to a
        SQLite database in WAL mode so that an external dashboard can read it while training is
        still in progress. \\
    FR-11 & Evaluation dashboard     & A headless evaluation mode and a Matplotlib dashboard
        mode shall be available, the latter showing LIDAR polar plot, speed gauge, lateral
        offset, shield status and intervention count. \\
    FR-12 & Lane-change awareness    & Both the reward shaper and the shield shall detect the
        condition \emph{lane change is permitted by road markings} and relax penalties that
        would otherwise block the manoeuvre. \\
    FR-13 & Reproducibility          & Training and evaluation shall be seeded independently
        (\(\mathrm{seed}=42\) and \(\mathrm{seed}=100\) respectively) and the seeds shall be
        exposed through command-line flags. \\
    \bottomrule
  \end{tabular}
\end{table}

\section{Non-functional requirements}\label{sec:req-nonfunctional}

Non-functional requirements describe qualities of the system: how fast, how portable, how
maintainable it must be. They are grouped below by ISO/IEC~25010 quality dimension.

\begin{table}[H]
  \centering
  \caption{Non-functional requirements of CALA-SafeRL.}
  \label{tab:nonfunctional-requirements}
  \small
  \begin{tabular}{@{}llp{7.8cm}@{}}
    \toprule
    \textbf{ID} & \textbf{Dimension} & \textbf{Description} \\
    \midrule
    NFR-01 & Performance             & The simulator shall run in synchronous mode at
        \(20\,\mathrm{Hz}\) (\(\Delta t = 0.05\,\mathrm{s}\)); all wrappers shall introduce an
        overhead small enough to sustain this rate on a single GPU workstation. \\
    NFR-02 & Determinism             & Given a fixed seed and a fixed hardware/software
        configuration, two training runs shall produce the same episode outcomes up to
        floating-point non-associativity. \\
    NFR-03 & Modularity              & The wrapper chain shall obey the Open/Closed principle:
        adding a new shield or a new reward component shall require no modifications to the
        existing layers. \\
    NFR-04 & Maintainability         & Python source files shall pass \texttt{ruff check} with
        the project's default configuration; public classes shall carry short docstrings. \\
    NFR-05 & Portability             & The framework shall run on Linux and Windows and shall
        automatically select the best PyTorch device among CUDA, Apple MPS and CPU. \\
    NFR-06 & Observability           & Every component that modifies actions or rewards shall
        expose an auxiliary \texttt{info} dictionary key (\texttt{executed\_action},
        \texttt{shield\_activated}, \texttt{risk\_level}, reward sub-components) so that
        training behaviour can be diagnosed post-hoc. \\
    NFR-07 & Extensibility           & The live-metrics schema shall admit new keys without
        breaking existing dashboards; both the SQLite logger and the Streamlit reader shall
        treat unknown metrics as first-class citizens. \\
    NFR-08 & Testability             & The reward shaper shall ship with deterministic unit
        tests that do not require a running CARLA server. \\
    NFR-09 & Documentation           & The project shall include a \texttt{CLAUDE.md} file that
        documents the wrapper chain, the observation space layout and the commands to reproduce
        training and evaluation. \\
    \bottomrule
  \end{tabular}
\end{table}

\section{Technical constraints}\label{sec:req-constraints}

The requirements above must be satisfied inside the technical envelope described in
Table~\ref{tab:technical-constraints}. These constraints are not design decisions but external
boundary conditions: the CARLA version is fixed by the research group's infrastructure, the
hardware is the author's workstation, and the timeline is set by the Bachelor's thesis calendar.

\begin{table}[H]
  \centering
  \caption{Technical constraints of the project.}
  \label{tab:technical-constraints}
  \small
  \begin{tabular}{@{}lp{10cm}@{}}
    \toprule
    \textbf{Category} & \textbf{Constraint} \\
    \midrule
    Simulator   & CARLA 0.9.16 running on Windows with the \texttt{-RenderOffScreen} flag
                  and fixed port \(2000\). \\
    Language    & Python \(\ge 3.10\). \\
    Core libs   & PyTorch (GPU when available), Gymnasium, NumPy, Matplotlib, Streamlit,
                  SQLite. Licences shall be compatible with open academic
                  distribution (MIT, BSD-3, Apache-2.0). \\
    Hardware    & Single desktop workstation; GPU with at least \(6\,\mathrm{GB}\) of VRAM is
                  required to run CARLA and the PyTorch policy concurrently. Training on CPU is
                  supported but orders of magnitude slower. \\
    Maps        & The main benchmark shall be Town04. Additional towns (Town01\textendash Town07,
                  Town10HD) may be used for qualitative evaluation but are not guaranteed to
                  converge with the default hyperparameters. \\
    Timeline    & The project shall be delivered within the Bachelor's thesis calendar of the
                  Universidad de Deusto. This precludes large hyperparameter sweeps that would
                  require multiple weeks of simulator time per configuration. \\
    Safety      & The framework shall only run inside simulation; deployment on a physical
                  vehicle is explicitly out of scope. \\
    \bottomrule
  \end{tabular}
\end{table}

\section{Acceptance criteria}\label{sec:req-acceptance}

Acceptance criteria translate the previous requirements into measurable conditions over the
metrics produced by \texttt{main\_train.py} and \texttt{main\_eval.py}. They are used in
Chapter~\ref{ch:experimentation} as the reference against which the three shielding
configurations are compared. The baseline corresponds to the unconstrained PPO agent
(\texttt{--shield\_type none}).

\begin{table}[H]
  \centering
  \caption{Acceptance criteria derived from the functional and non-functional requirements.
           Thresholds are placeholders to be replaced by the empirical numbers obtained in
           Chapter~\ref{ch:experimentation}.}
  \label{tab:acceptance-criteria}
  \small
  \begin{tabular}{@{}llp{7.5cm}@{}}
    \toprule
    \textbf{ID} & \textbf{Metric} & \textbf{Acceptance threshold} \\
    \midrule
    AC-01 & Reward convergence     & The 100-episode rolling average reward of the shielded
        agent shall be no worse than the baseline after an equal amount of training. \\
    AC-02 & Crash rate             & With the adaptive shield, the 100-episode rolling
        \texttt{crash\_rate} shall drop by at least a factor of \(\alpha=4\) with respect to the
        unshielded probe (empirically: \(16\%\to 2\%\), a factor of \(8\times\); the threshold
        is set conservatively at \(4\times\)). \\
    AC-03 & Off-road rate          & Under the adaptive shield, the 100-episode rolling
        \texttt{offroad\_rate} shall remain below \(\beta=0.15\) (\(15\%\)).
        Empirically: \(0.0\%\) at training end; \(12\%\) in the \(50\)-episode evaluation. \\
    AC-04 & Success rate           & On \(50\) evaluation episodes with seed \(100\), the agent
        shall reach the \texttt{success\_distance} target in at least \(\gamma=0.60\)
        (\(60\%\)) of them. Empirically: \(70\%\) with adaptive shield. \\
    AC-05 & Shield interference    & The fraction of all steps in which the adaptive shield
        activates shall not exceed \(\tau=0.40\) (\(40\%\)). The \emph{safe}, \emph{warning}
        and \emph{critical} risk-level fractions shall sum to \(1\). Empirically: shield
        activates in \(38.7\%\) of steps; \(65.8\%\) safe, \(25.6\%\) warning, \(8.6\%\)
        critical. \\
    AC-06 & Speed compliance       & At least \(90\%\) of steps with a valid speed limit shall
        satisfy \(\text{speed\_kmh} \le 1.05\cdot\text{speed\_limit\_kmh}\). \\
    AC-07 & Curriculum progression & The curriculum manager shall advance from stage 0 to at
        least stage 3 during a standard \(2\,500\)-episode run in all shielded
        configurations. \\
    AC-08 & Dashboard liveness     & The Streamlit dashboard shall refresh training metrics
        with an end-to-end latency below \(2\,\mathrm{s}\) from the moment the logger writes a
        row to the SQLite database. \\
    AC-09 & Reproducibility        & Two independent runs with the same seed and configuration
        shall produce identical reward curves in the first \(50\) episodes (before floating-point
        divergence becomes significant). \\
    \bottomrule
  \end{tabular}
\end{table}

The criteria with placeholder thresholds (\(\alpha, \beta, \gamma, \tau\)) are explicitly left to
be fixed once the experiments of Chapter~\ref{ch:experimentation} have been executed. This keeps
the requirements traceable while avoiding the pitfall of committing to numeric targets before the
empirical evidence is available.

\clearpage
